\section*{Bloque D. Ampliación: Teoría y Aplicaciones Prácticas (70\% / 30\%)}

\noindent\textit{Preguntas 91 a 180. Ampliación del banco oficial con 63 preguntas teóricas y conceptuales (91 a 153) y 27 aplicaciones cuantitativas prácticas (154 a 180).}

\subsection*{Parte I. Preguntas teóricas y conceptuales (91 a 153)}

\q{Cuando una tasa nominal anual $r > 0$ se capitaliza un número $m$ de veces al año y se toma el límite conforme $m \to \infty$ (capitalización continua), la Tasa Efectiva Anual (TEA) converge exactamente a:}
\begin{opts}
    \item $TEA = r \cdot \ln(1 + r)$, lo cual produce un rendimiento menor que la capitalización diaria.
    \item $TEA = e^{r} - 1$, representando el límite superior efectivo alcanzable para esa tasa nominal.
    \item $TEA = (1 + r)^e - 1$, creciendo sin cota superior conforme aumenta la frecuencia de capitalización.
    \item $TEA = \frac{r}{1 - e^{-r}}$, valor que coincide exactamente con el interés simple anual.
\end{opts}

\q{En un contrato de deuda pactado a tasa de interés nominal fija sin cláusula de indexación, si la inflación efectiva (ex-post) resulta sustancialmente mayor que la inflación esperada al momento de firmar el contrato, entonces:}
\begin{opts}
    \item La tasa real ex-post aumenta por encima de la esperada, redistribuyendo riqueza de los deudores hacia los acreedores.
    \item Tanto la tasa real ex-ante como la tasa real ex-post permanecen inalteradas según la ecuación exacta de Fisher.
    \item La tasa real ex-post cae por debajo de la esperada, beneficiando a los deudores en perjuicio de los acreedores.
    \item El valor presente real de los pagos futuros de la deuda aumenta proporcionalmente a la sorpresa inflacionaria.
\end{opts}

\q{En el modelo de valoración de acciones de Gordon-Shapiro ($P_0 = \frac{D_1}{k_e - g}$), si una empresa retiene una fracción $b > 0$ de sus utilidades para reinvertirlas en proyectos cuya rentabilidad es exactamente igual al costo del patrimonio ($ROE = k_e$), entonces el Valor Presente de las Oportunidades de Crecimiento (PVGO):}
\begin{opts}
    \item Es estrictamente positivo porque la tasa de crecimiento de dividendos $g = ROE \cdot b$ es mayor que cero.
    \item Tiende al infinito conforme la tasa de retención $b$ se aproxima al $100\%$ de la utilidad neta anual.
    \item Es estrictamente negativo porque la retención reduce el dividendo inmediato sin elevar la rentabilidad exigida.
    \item Es igual a cero, pues retener utilidades no crea valor adicional frente a repartirlas.
\end{opts}

\q{Un crédito de inversión para infraestructura de agua potable contempla una anualidad vencida de $n$ cuotas iguales de monto $C$ que comienzan a pagarse al final del año $k + 1$ (es decir, cuenta con $k$ años de periodo de gracia sin pagos). Su valor presente en $t = 0$ se obtiene:}
\begin{opts}
    \item Hallando el valor presente de la anualidad en $t = k$ y descontándolo por $(1 + r)^k$ a $t = 0$.
    \item Sumando $k$ cuotas adicionales al plazo $n$ y descontando el total con el factor de anualidad anticipada.
    \item Multiplicando el valor presente de una anualidad inmediata por $(1 + r)^k$ para compensar el diferimiento.
    \item Dividiendo directamente la cuota fija $C$ entre la tasa $r$ sin considerar el número de años de gracia $k$.
\end{opts}

\q{Según la teoría de la preferencia por la liquidez aplicada a la estructura temporal de tasas de interés (curva de rendimiento), incluso si el mercado espera que las tasas de interés de corto plazo permanezcan constantes en el futuro, la curva de rendimiento tenderá a presentar:}
\begin{opts}
    \item Pendiente negativa (invertida), debido a que los inversionistas exigen un descuento por inmovilizar capital a largo plazo.
    \item Una forma perfectamente horizontal, dado que las tasas *forward* coinciden exactamente con las tasas *spot* esperadas.
    \item Pendiente positiva (ascendente), porque los bonos de mayor plazo incorporan una prima por liquidez creciente.
    \item Una forma cóncava con tasas nominales negativas en los plazos intermedios por efecto de la segmentación.
\end{opts}

\q{En un mercado de renta fija sin fricciones ni oportunidades de arbitraje, la tasa *forward* implícita $f_{1,2}$ para el segundo año se vincula con las tasas *spot* anuales $s_1$ (a un año) y $s_2$ (a dos años) mediante la condición:}
\begin{opts}
    \item $(1 + f_{1,2})^2 = (1 + s_1)(1 + s_2)$, representando el promedio geométrico exacto de ambas tasas *spot*.
    \item $(1 + s_2) = (1 + s_1)(1 + f_{1,2})$, sumando simplemente los intereses compuestos de ambos horizontes.
    \item $f_{1,2} = s_2 - s_1$, lo que implica que la tasa *forward* nunca puede superar la tasa *spot* de mayor plazo.
    \item $(1 + s_2)^2 = (1 + s_1)(1 + f_{1,2})$, igualando el retorno de invertir a 2 años con reinvertir a 1 año.
\end{opts}

\q{Al comparar el servicio de la deuda de un préstamo bajo el sistema francés (cuota constante) frente al sistema alemán (amortización de capital constante) para el mismo capital, plazo y tasa de interés:}
\begin{opts}
    \item En el sistema francés las cuotas totales son constantes, con intereses decrecientes y amortización creciente.
    \item En el sistema francés la amortización de capital es constante y las cuotas totales decrecen cada periodo.
    \item En el sistema alemán las cuotas totales son crecientes porque los intereses se capitalizan hasta el vencimiento.
    \item Ambos sistemas generan exactamente el mismo total acumulado de intereses pagados a lo largo de la vida del crédito.
\end{opts}

\q{Para una Tasa Nominal Anual (TNA) estrictamente positiva y fija, incrementar la frecuencia de capitalización $m$ (por ejemplo, pasando de capitalización semestral a mensual y luego a diaria) ocasiona que la Tasa Efectiva Anual (TEA):}
\begin{opts}
    \item Disminuya progresivamente acercándose a la TNA debido a la reducción de la tasa periódica $TNA/m$.
    \item Aumente de manera estrictamente creciente pero acotada superiormente por $e^{TNA} - 1$.
    \item Permanezca completamente invariante porque la TNA define por sí sola el costo financiero efectivo anual.
    \item Crezca linealmente sin límite conforme aumenta el número de periodos de capitalización $m$.
\end{opts}

\q{¿Por qué una empresa de servicios públicos puede reportar una utilidad neta contable positiva en su Estado de Resultados y, simultáneamente, sufrir un déficit severo de efectivo que compromete su liquidez?}
\begin{opts}
    \item Porque las inversiones en infraestructura de largo plazo se cargan íntegramente como gasto operativo en el año cero.
    \item Porque el Estado de Resultados registra únicamente entradas en efectivo y excluye los gastos de depreciación contable.
    \item Porque los gastos financieros por intereses no se deducen en el cálculo de la utilidad neta del ejercicio.
    \item Porque sigue el devengado y omite salidas de caja como CAPEX o alzas en cuentas por cobrar.
\end{opts}

\q{En finanzas corporativas, la regla de consistencia entre el flujo de caja proyectado y la tasa de descuento establece que el Flujo de Caja Libre de la Firma (FCL / FCFF) y el Flujo de Caja del Accionista (FCA / FCFE) deben descontarse, respectivamente, a:}
\begin{opts}
    \item El Costo Promedio Ponderado de Capital (WACC) para el FCL y el costo del patrimonio ($k_e$) para el FCA.
    \item El costo del patrimonio ($k_e$) para el FCL y el costo de la deuda antes de impuestos ($k_d$) para el FCA.
    \item La tasa libre de riesgo ($r_f$) para el FCL y el WACC sin escudo fiscal para el FCA.
    \item El WACC para ambos flujos, ajustando el FCA únicamente por la tasa de inflación esperada del país.
\end{opts}

\q{Al construir el Flujo de Caja Libre (FCL) de un proyecto de inversión de horizonte finito, un incremento positivo del Capital de Trabajo Neto ($\Delta CTN_t > 0$) durante la fase operativa y su liquidación al término del proyecto se registran, respectivamente, como:}
\begin{opts}
    \item Una entrada de efectivo durante la operación y una salida de efectivo al cierre del horizonte de evaluación.
    \item Una salida de efectivo durante la operación y una recuperación (entrada) de efectivo en el último año.
    \item Un gasto deducible del impuesto a la renta en la operación sin ningún impacto en el último periodo.
    \item Partidas estrictamente contables que no modifican el flujo de caja libre ni el Valor Actual Neto.
\end{opts}

\q{Aunque la depreciación contable no constituye una salida real de dinero en el periodo en que se registra, su inclusión en el flujo de caja de un proyecto privado sometido a impuesto a la renta ($T > 0$) es financieramente relevante porque:}
\begin{opts}
    \item Aumenta el Costo Promedio Ponderado de Capital (WACC) al elevar el riesgo sistemático del negocio.
    \item Incrementa directamente la utilidad neta contable y el monto de impuestos a pagar a la autoridad tributaria.
    \item Reduce la base imponible y crea un escudo fiscal efectivo igual a $T \times \text{Dep}$.
    \item Sustituye completamente el desembolso inicial de inversión (CAPEX) en la fórmula del Valor Actual Neto.
\end{opts}

\q{El Valor Económico Agregado ($EVA = NOPAT - WACC \times \text{Capital Invertido}$) permite evaluar la creación de valor corporativo en un periodo. Una empresa destruye valor económico ($EVA < 0$) aun con utilidad operativa positiva cuando:}
\begin{opts}
    \item El retorno sobre el capital invertido ($ROIC$) es inferior al Costo Promedio Ponderado de Capital ($WACC$).
    \item El retorno sobre el capital invertido ($ROIC$) supera estrictamente al Costo Promedio Ponderado de Capital ($WACC$).
    \item La razón de endeudamiento $D/E$ es igual a cero y la empresa se financia íntegramente con capital propio.
    \item La depreciación del ejercicio es mayor que el pago de dividendos en efectivo a los accionistas.
\end{opts}

\q{En la descomposición ampliada de cinco factores de la identidad DuPont, el margen neto ($\text{UN}/\text{Ventas}$) se desglosa en tres componentes multiplicativos que permiten separar el desempeño operativo puro del efecto financiero y tributario:}
\begin{opts}
    \item Margen bruto, razón corriente de liquidez y periodo promedio de cobranza comercial.
    \item Margen operativo ($EBIT/V$), carga financiera ($EBT/EBIT$) y tributaria ($UN/EBT$).
    \item Rotación de inventarios, cobertura de intereses y tasa de retención de utilidades.
    \item Margen EBITDA, relación deuda/activos y costo promedio ponderado de la deuda bancaria.
\end{opts}

\q{El apalancamiento financiero resulta favorable para los accionistas (es decir, incrementa la rentabilidad sobre el patrimonio, $ROE$, a medida que sube la razón deuda/patrimonio $D/E$) siempre que:}
\begin{opts}
    \item La empresa mantenga pérdidas operativas recurrentes que impidan aprovechar el escudo fiscal de los intereses.
    \item El costo de la deuda antes de impuestos ($k_d$) supere ampliamente el rendimiento operativo de los activos ($ROIC$).
    \item El rendimiento operativo neto ($ROIC$) supere el costo efectivo de la deuda $k_d(1 - T)$.
    \item La rotación de activos totales disminuya proporcionalmente al incremento del pasivo financiero.
\end{opts}

\q{¿Qué diferencia existe entre el Grado de Apalancamiento Operativo (GAO) y el Grado de Apalancamiento Financiero (GAF) de una empresa?}
\begin{opts}
    \item El GAO depende exclusivamente del nivel de endeudamiento bancario, mientras que el GAF depende de los costos variables unitarios.
    \item El GAO mide la sensibilidad de las ventas ante la tasa de interés; el GAF mide la proporción de activos financiados con capital propio.
    \item Ambos miden exactamente el mismo riesgo de mercado sistemático y se suman para obtener el Grado de Apalancamiento Total.
    \item GAO mide sensibilidad del $EBIT$ a las ventas (costos fijos); GAF mide sensibilidad de $UPA$ al $EBIT$ (intereses).
\end{opts}

\q{Una EPS evalúa construir una nueva planta de tratamiento sobre un terreno de su propiedad adquirido hace diez años, el cual podría venderse hoy en el mercado en $S/\ 2$ millones. En el flujo de caja incremental del proyecto, dicho terreno debe incorporarse como:}
\begin{opts}
    \item Un costo hundido igual a cero, dado que el terreno ya pertenece a la empresa y no implica desembolso de caja hoy.
    \item Un costo de oportunidad de $S/\ 2$ millones (neto de impuestos sobre ganancia de capital) en el año $t = 0$.
    \item Su valor histórico de compra de hace diez años registrado en el balance contable sin actualizarlo a precios de mercado.
    \item Un ingreso extraordinario de $S/\ 2$ millones en el año $t = 0$ por tratarse de un aporte en especie.
\end{opts}

\q{Según la regla de los signos de Descartes aplicada al polinomio del VAN de un proyecto con flujos no convencionales (por ejemplo, inversión inicial negativa, flujos operativos positivos y un fuerte costo de cierre ambiental negativo en el último año):}
\begin{opts}
    \item El proyecto siempre posee una única Tasa Interna de Retorno real estrictamente positiva.
    \item El criterio del Valor Actual Neto (VAN) deja de ser válido y debe reemplazarse obligatoriamente por el Payback.
    \item Puede existir hasta una TIR real positiva por cada cambio de signo en la secuencia de flujos.
    \item La Tasa Interna de Retorno Modificada (TIRM) también produce múltiples soluciones matemáticas irreconciliables.
\end{opts}

\q{La Tasa de Intersección de Fisher ($r_F$) entre dos proyectos mutuamente excluyentes $A$ y $B$ se define formalmente como la tasa de descuento en la cual $VAN_A(r_F) = VAN_B(r_F)$. Esta tasa coincide matemáticamente con:}
\begin{opts}
    \item El promedio aritmético simple de las Tasas Internas de Retorno individuales de ambos proyectos.
    \item La tasa de descuento a la cual el Índice de Rentabilidad de ambos proyectos es exactamente igual a cero.
    \item El Costo Promedio Ponderado de Capital (WACC) mínimo que hace positivos ambos Valores Actuales Netos.
    \item La TIR del proyecto incremental obtenido al restar sus flujos ($FC_A - FC_B$).
\end{opts}

\q{Cuando una empresa enfrenta un racionamiento duro de capital en un solo periodo (presupuesto máximo fijo $B_0$) y dispone de múltiples proyectos independientes y divisibles con $VAN > 0$, la regla óptima de asignación consiste en ordenar los proyectos de mayor a menor según:}
\begin{opts}
    \item Su Índice de Rentabilidad ($IR = VP/I_0$), maximizando el VAN por sol invertido.
    \item Su Periodo de Recuperación Simple (Payback) para recuperar el efectivo lo más rápido posible.
    \item Su Valor Actual Neto absoluto individual, sin considerar el tamaño de la inversión inicial requerida.
    \item Su utilidad contable promedio anual dividida entre la vida útil contable del activo.
\end{opts}

\q{Para comparar dos proyectos mutuamente excluyentes y repetibles con vidas útiles distintas (por ejemplo, una bomba de 4 años frente a otra de 6 años), ¿por qué es incorrecto elegir simplemente el de menor valor presente de costos (VAC) de un solo ciclo?}
\begin{opts}
    \item Porque el VAC de un solo ciclo no incorpora la tasa de descuento del costo de oportunidad del capital.
    \item Porque la bomba de menor vida útil siempre tiene un costo anual equivalente menor independientemente de su inversión.
    \item Porque compara horizontes temporales desiguales; debe usarse el Costo Anual Equivalente (CAE) o cadena de reemplazos.
    \item Porque el criterio del VAN solo es aplicable a proyectos de vida útil infinita o perpetuidades crecientes.
\end{opts}

\q{Al evaluar el momento óptimo de reemplazo de una maquinaria existente por una nueva de igual función, la decisión económicamente óptima consiste en reemplazar la máquina antigua tan pronto como:}
\begin{opts}
    \item Su valor en libros contable llegue exactamente a cero tras completar su depreciación tributaria.
    \item Los gastos de mantenimiento acumulados desde el año cero igualen el precio de compra original del equipo.
    \item Su valor de rescate en el mercado secundario sea estrictamente menor que su costo de adquisición histórico.
    \item El costo marginal anual de mantenerla supere el Costo Anual Equivalente (CAE) de la nueva.
\end{opts}

\q{El Periodo de Recuperación Descontado (*Discounted Payback*) mejora al Payback simple al incorporar el valor del dinero en el tiempo; sin embargo, mantiene una limitación crítica que le impide sustituir al VAN:}
\begin{opts}
    \item Ignora por completo todos los flujos de caja que ocurren después de la fecha de corte de recuperación.
    \item Exige conocer una tasa de reinversión diferente del costo de oportunidad del capital para cada periodo.
    \item Produce resultados negativos siempre que los flujos de caja sean constantes durante toda la vida del proyecto.
    \item No puede calcularse cuando el proyecto presenta una única inversión inicial seguida de flujos positivos.
\end{opts}

\q{En la evaluación de proyectos bajo incertidumbre, la principal ventaja de la Simulación de Monte Carlo frente al análisis de sensibilidad univariado tradicional radica en que:}
\begin{opts}
    \item Elimina por completo la necesidad de estimar el Costo Promedio Ponderado de Capital (WACC) del proyecto.
    \item Varía simultáneamente variables inciertas según sus distribuciones y obtiene la distribución del VAN.
    \item Garantiza que el Valor Actual Neto esperado del proyecto sea siempre estrictamente positivo y libre de riesgo.
    \item Sustituye los flujos de caja proyectados por valores históricos contables determinísticos de los últimos cinco años.
\end{opts}

\q{Desde el enfoque de Opciones Reales, ¿por qué el método del VAN estático tradicional tiende a subestimar el valor de grandes proyectos de infraestructura con alta incertidumbre y fases modulares?}
\begin{opts}
    \item Porque el VAN estático sobreestima los costos fijos operativos cuando la volatilidad de la demanda es elevada.
    \item Porque el VAN estático utiliza tasas de descuento nominales aplicadas sobre flujos expresados en moneda constante.
    \item Porque la metodología de opciones reales elimina el costo de inversión inicial del cálculo del valor presente.
    \item Porque asume una decisión rígida de 'ahora o nunca' e ignora la flexibilidad gerencial futura.
\end{opts}

\q{En la teoría moderna de portafolios de Markowitz, al combinar dos activos riesgosos con retornos esperados y varianzas positivas, la diversificación logra reducir la desviación estándar del portafolio por debajo del promedio ponderado de sus desviaciones individuales siempre que:}
\begin{opts}
    \item La correlación lineal entre ambos activos sea estrictamente menor que $+1$ ($\rho_{12} < 1$).
    \item El coeficiente de correlación lineal entre ambos activos sea estrictamente igual a $+1$ ($\rho_{12} = 1$).
    \item Ambos activos presenten exactamente la misma tasa de retorno esperado y la misma varianza individual.
    \item Se invierta más del $50\%$ del portafolio en el activo libre de riesgo antes de combinarlos.
\end{opts}

\q{Según el Teorema de Separación de Dos Fondos de Tobin (cuando existe un activo libre de riesgo y mercados sin fricciones), todos los inversionistas aversos al riesgo conforman su portafolio óptimo:}
\begin{opts}
    \item Eligiendo distintas combinaciones internas de acciones individuales según su grado personal de aversión al riesgo.
    \item Combinando el activo libre de riesgo con el mismo portafolio tangente según su aversión al riesgo.
    \item Invirtiendo el $100\%$ de su riqueza en el portafolio de mínima varianza global sin utilizar el activo libre de riesgo.
    \item Seleccionando exclusivamente activos cuya beta individual sea estrictamente negativa para anular el riesgo.
\end{opts}

\q{La distinción clave entre la Línea del Mercado de Capitales (CML) y la Línea del Mercado de Valores (SML) en el modelo CAPM es que:}
\begin{opts}
    \item La pendiente de la CML es la prima de riesgo de mercado ($E[R_m] - r_f$), mientras que la de la SML es el ratio de Sharpe.
    \item La CML aplica a cualquier acción individual, mientras que la SML solo aplica a bonos soberanos libres de riesgo.
    \item CML usa $\sigma_p$ solo para portafolios eficientes; SML usa $\beta_i$ para todo activo.
    \item La SML exige que el activo esté perfectamente diversificado, en tanto la CML incluye activos ineficientes.
\end{opts}

\q{Si un analista estima que una acción ofrece un rendimiento esperado que se ubica estrictamente por encima de la Línea del Mercado de Valores (SML), lo que equivale a un Alfa de Jensen positivo ($\alpha_i > 0$), en equilibrio esto indica que el activo está:}
\begin{opts}
    \item Subvaluado frente a su riesgo sistemático, ofreciendo oportunidad de compra.
    \item Sobrevaluado respecto de su riesgo sistemático, por lo que conviene venderlo en corto.
    \item Correctamente valorado en un mercado eficiente sin posibilidad de obtener retornos ajustados por riesgo.
    \item Afectado por un riesgo no sistemático excesivo que reduce su rendimiento esperado por debajo de $r_f$.
\end{opts}

\q{Al evaluar el desempeño ajustado por riesgo de un fondo de inversión, ¿en qué caso resulta preferible utilizar el Índice de Treynor en lugar del Ratio de Sharpe?}
\begin{opts}
    \item Cuando el fondo representa la totalidad de la riqueza del inversionista y el riesgo relevante es la desviación estándar total.
    \item Cuando el fondo integra una cartera diversificada y solo importa su riesgo sistemático ($\beta$).
    \item Cuando el fondo no posee exposición alguna al mercado accionario y su beta es exactamente igual a cero.
    \item Cuando se desea medir exclusivamente la rentabilidad bruta sin descontar la tasa libre de riesgo $r_f$.
\end{opts}

\q{Para estimar el beta del patrimonio de una empresa de agua que no cotiza en bolsa, el regulador suele aplicar el método de comparables (*pure-play*) con la ecuación de Hamada. El procedimiento correcto consiste en:}
\begin{opts}
    \item Promediar directamente los betas apalancados de empresas comparables sin ajustar por sus diferencias de endeudamiento ni tasas impositivas.
    \item Multiplicar el beta del mercado bursátil local por la razón de liquidez corriente de la empresa regulada.
    \item Desapalancar betas comparables con su $D/E$ y $T$, y reapalancar $\beta_U$ con la $D/E$ de la EPS.
    \item Estimar el beta mediante una regresión de la utilidad contable de la EPS contra la inflación anual del país.
\end{opts}

\q{De acuerdo con la Proposición II de Modigliani-Miller en un mundo sin impuestos ni costos de quiebra, aunque la deuda sea más barata que el patrimonio ($k_d < k_e$), aumentar el endeudamiento $D/E$ no reduce el WACC de la empresa porque:}
\begin{opts}
    \item El valor de mercado de los activos cae en la misma proporción en que aumenta el pasivo financiero.
    \item El costo de la deuda $k_d$ aumenta instantáneamente hasta igualar el costo del patrimonio inicial $k_u$.
    \item Los impuestos corporativos absorben íntegramente el ahorro financiero generado por los intereses pagados.
    \item El mayor apalancamiento eleva $k_e$ exactamente lo necesario para compensar la deuda barata.
\end{opts}

\q{Según la teoría estática del intercambio (*Trade-off Theory*) de la estructura de capital, una empresa alcanza su nivel óptimo de endeudamiento $D^*$ cuando:}
\begin{opts}
    \item Se financia al $100\%$ con deuda para maximizar el valor presente del escudo fiscal corporativo.
    \item El beneficio marginal del escudo fiscal iguala el costo marginal esperado de quiebra.
    \item Se financia al $100\%$ con capital propio para eliminar completamente la probabilidad de incumplimiento financiero.
    \item El costo del patrimonio $k_e$ se vuelve estrictamente menor que la tasa de interés de la deuda libre de riesgo.
\end{opts}

\q{La teoría del orden jerárquico (*Pecking Order Theory*) de Myers y Majluf explica que, debido a la asimetría de información entre gerentes e inversionistas externos, las empresas prefieren financiar nuevos proyectos siguiendo el orden:}
\begin{opts}
    \item Primero emisión de nuevas acciones comunes, segundo deuda bancaria y finalmente utilidades retenidas.
    \item Primero deuda convertible subordinada, segundo emisión de acciones preferentes y finalmente flujo de caja operativo.
    \item Primero recursos internos (utilidades retenidas), segundo emisión de deuda y como último recurso emisión de acciones.
    \item Indistintamente deuda o acciones, pues en presencia de asimetría informativa la estructura de capital es irrelevante.
\end{opts}

\q{En una empresa altamente endeudada y cercana a la insolvencia, el problema de agencia conocido como 'sustitución de activos' (*asset substitution* o toma excesiva de riesgo) ocurre cuando los accionistas tienen incentivos a:}
\begin{opts}
    \item Rechazar proyectos seguros con VAN positivo porque todos sus beneficios irían a los acreedores.
    \item Mantener todo el efectivo en bonos del Tesoro libre de riesgo para proteger la solvencia de los bonistas.
    \item Emitir acciones nuevas para prepagar anticipadamente toda la deuda bancaria antes del vencimiento.
    \item Elegir proyectos muy riesgosos porque capturan ganancias altas y trasladan pérdidas a acreedores.
\end{opts}

\q{De acuerdo con el teorema de irrelevancia de la política de dividendos de Miller y Modigliani (1961), en mercados perfectos sin impuestos ni costos de transacción y dada la política de inversiones de la empresa:}
\begin{opts}
    \item La política de dividendos es irrelevante, pues los accionistas pueden crear dividendos caseros.
    \item Pagar un dividendo mayor siempre aumenta el precio de la acción porque los inversionistas prefieren efectivo inmediato.
    \item Retener el $100\%$ de las utilidades maximiza el valor de la empresa independientemente del VAN de los proyectos.
    \item El precio de la acción sube exactamente en el monto del dividendo el día ex-dividendo.
\end{opts}

\q{¿En qué situación resulta metodológicamente superior utilizar el método del Valor Presente Ajustado (APV) en lugar del enfoque tradicional de descontar flujos al WACC constante?}
\begin{opts}
    \item Cuando la empresa mantiene una razón deuda/valor ($D/V$) estrictamente constante durante todos los años del proyecto.
    \item Cuando el proyecto se financia íntegramente con capital propio y no existen efectos tributarios ni subsidios financieros.
    \item Cuando la razón deuda/valor ($D/V$) cambia cada año (como en *Project Finance* o deuda fija).
    \item Cuando se desconoce por completo la tasa impositiva corporativa y el calendario de pagos de la deuda.
\end{opts}

\q{Un bono corporativo con tasa cupón fija del $6\%$ anual cotiza actualmente en el mercado secundario con un rendimiento al vencimiento (YTM) del $7.5\%$. En consecuencia, este bono se negocia:}
\begin{opts}
    \item Sobre la par (con prima), a un precio estrictamente superior a su valor nominal.
    \item A un precio igual al valor presente de sus cupones sin incluir la devolución del principal.
    \item Exactamente a la par, dado que el valor nominal no cambia con las variaciones del YTM.
    \item Bajo la par (con descuento), a un precio estrictamente inferior a su valor nominal.
\end{opts}

\q{Respecto de las propiedades de la Duración de Macaulay ($D$) de un bono convencional con tasa cupón positiva y rendimiento al vencimiento $y > 0$, es correcto afirmar que, manteniendo lo demás constante:}
\begin{opts}
    \item Es menor que $n$, cae al subir la tasa cupón y coincide con $n$ en un bono cupón cero.
    \item La duración es estrictamente mayor que el plazo al vencimiento $n$ y aumenta conforme sube la tasa cupón.
    \item La duración es independiente del nivel de la tasa de rendimiento $y$ y solo depende del valor nominal del bono.
    \item La duración disminuye proporcionalmente conforme se alarga el plazo al vencimiento de un bono cupón cero.
\end{opts}

\q{Debido a la convexidad positiva de la curva precio-rendimiento de un bono tradicional (sin opciones implícitas), utilizar exclusivamente la aproximación lineal basada en la Duración Modificada ($\frac{\Delta P}{P} \approx -D_{mod}\Delta y$) ante cambios grandes en la tasa de interés:}
\begin{opts}
    \item Sobreestima el precio real tanto cuando las tasas suben como cuando las tasas bajan.
    \item Subestima el precio real tanto cuando las tasas suben como cuando las tasas bajan.
    \item Subestima el precio cuando las tasas suben, pero sobreestima el precio cuando las tasas bajan.
    \item Predice el precio exacto sin error alguno sin importar la magnitud del cambio en el rendimiento.
\end{opts}

\q{En la gestión de carteras de renta fija (por ejemplo, un fondo de reserva de una EPS para pagar una obligación única dentro de $5$ años), la estrategia clásica de inmunización contra pequeños desplazamientos paralelos de la tasa de interés consiste en:}
\begin{opts}
    \item Invertir exclusivamente en bonos cuyo plazo al vencimiento sea superior a $20$ años.
    \item Mantener todo el fondo en efectivo sin intereses para anular el riesgo de reinversión.
    \item Invertir únicamente en bonos con la mayor tasa cupón disponible sin considerar su duración.
    \item Formar una cartera cuya Duración de Macaulay iguale el horizonte de $5$ años.
\end{opts}

\q{La relación de Paridad Put-Call para opciones europeas sobre una acción que no paga dividendos ($C + \frac{K}{(1+r)^T} = P + S_0$) se demuestra mediante el principio de ausencia de arbitraje porque:}
\begin{opts}
    \item Una *call* más un bono de nominal $K$ replican los pagos de una *put* más la acción.
    \item Las primas de una opción *call* y una opción *put* con igual precio de ejercicio deben ser siempre numéricamente idénticas.
    \item El precio de ejercicio $K$ debe coincidir obligatoriamente con el precio spot actual $S_0$ en la fecha de emisión.
    \item Las opciones europeas pueden ejercerse en cualquier fecha anterior al vencimiento $T$ sin penalidad.
\end{opts}

\q{Según el modelo de valoración de opciones de Black-Scholes-Merton, un incremento en la volatilidad ($\sigma$) del activo subyacente —manteniendo constantes los demás parámetros— provoca que:}
\begin{opts}
    \item Aumente el valor de una opción *call* europea y disminuya el valor de una opción *put* europea.
    \item Aumenten simultáneamente tanto el precio de la *call* como el de la *put* europea.
    \item Disminuyan tanto el valor de la *call* como el de la *put* por mayor aversión al riesgo del mercado.
    \item No cambie el valor de ninguna opción si el precio actual coincide con el precio de ejercicio.
\end{opts}

\q{Una EPS que contrae un préstamo en dólares mientras sus ingresos tarifarios están en soles decide cubrir su riesgo cambiario. ¿Qué diferencia esencial existe entre contratar un *Forward* de compra de dólares y comprar una opción *Call* sobre dólares?}
\begin{opts}
    \item El *Forward* exige pagar una prima inicial no reembolsable elevada, mientras que la opción *Call* tiene costo inicial cero.
    \item La opción *Call* obliga a comprar dólares al tipo de cambio pactado aunque el dólar caiga en el mercado spot.
    \item El *Forward* fija el tipo de cambio obligatoriamente; la *Call* pone un techo pagando una prima.
    \item Ambos instrumentos eliminan completamente el riesgo crediticio y presentan exactamente el mismo perfil lineal de pagos.
\end{opts}

\q{En la taxonomía de Opciones Reales aplicadas a proyectos de infraestructura, la opción de diferir (esperar antes de invertir) y la opción de abandonar un proyecto vendiendo sus activos a un valor de salvamento son análogas, respectivamente, a:}
\begin{opts}
    \item Una *call* (ejercicio igual a inversión) y una *put* (ejercicio igual al salvamento).
    \item Una opción *put* sobre el proyecto y una opción *call* sobre el valor de salvamento.
    \item Un contrato *forward* sobre tarifas y un *swap* de tasas de interés de largo plazo.
    \item Dos opciones *put* europeas cuyo valor disminuye conforme aumenta la incertidumbre de la demanda.
\end{opts}

\q{En el marco tarifario de SUNASS para las empresas prestadoras de servicios de saneamiento (EPS), la determinación de la tarifa media de mediano plazo busca cumplir la condición financiera fundamental de que:}
\begin{opts}
    \item Los ingresos tarifarios anuales igualen únicamente los costos variables de energía y químicos sin remunerar el capital.
    \item El VP de ingresos tarifarios iguale el VP de costos eficientes de O\&M, CAPEX e impuestos al WACC.
    \item La Tasa Interna de Retorno (TIR) de la EPS sea exactamente igual a cero en todos los años del periodo regulatorio.
    \item Las tarifas cubran íntegramente las ineficiencias operativas históricas y las pérdidas por agua no facturada.
\end{opts}

\q{En la estimación del WACC regulatorio por parte de SUNASS para el sector saneamiento peruano, ¿por qué se utiliza una estructura de capital objetivo (referencial) y un beta de empresas comparables en lugar de los datos contables propios de cada EPS municipal?}
\begin{opts}
    \item Porque todas las EPS del Perú tienen exactamente el mismo nivel de activos y la misma razón deuda/patrimonio.
    \item Porque la ley prohíbe explícitamente utilizar la tasa de impuesto a la renta peruana en el cálculo del WACC.
    \item Porque las EPS no cotizan en bolsa y presentan estructuras contables de deuda distorsionadas.
    \item Porque el modelo CAPM solo puede aplicarse utilizando exclusivamente empresas del sector financiero bancario.
\end{opts}

\q{Al valorar la Base de Capital (Base de Activos Regulatorios) de una empresa de servicios públicos mediante el Valor Nuevo de Reemplazo (VNR) en lugar del costo histórico contable depreciado, una consecuencia directa es que:}
\begin{opts}
    \item Se elimina la necesidad de estimar la tasa de descuento WACC en el estudio tarifario quinquenal.
    \item Se garantiza que la tarifa disminuya a cero una vez que los activos antiguos cumplen su vida útil contable.
    \item Se remuneran obligatoriamente todos los activos ociosos o sobredimensionados construidos en décadas pasadas.
    \item Refleja el costo actual eficiente de reponer la infraestructura, evitando sesgos históricos.
\end{opts}

\q{Cuando una EPS recibe infraestructura de agua o saneamiento financiada mediante transferencias gratuitas o donaciones del Gobierno Nacional (por ejemplo, del MVCS/PNSU), el tratamiento regulatorio estándar en el estudio tarifario de SUNASS consiste en:}
\begin{opts}
    \item Incluir dichas obras donadas en la Base de Capital remunerada al WACC para generar utilidades extraordinarias a la EPS.
    \item Excluir el capital donado del WACC, pero incluir en tarifa los costos eficientes de O\&M y reposición.
    \item Excluir tanto la inversión como los costos futuros de operación y mantenimiento de las obras transferidas.
    \item Cobrar a los usuarios el $100\%$ del valor de la obra donada en el primer año del periodo regulatorio.
\end{opts}

\q{En la regulación económica de SUNASS, ¿cómo se articulan el Plan Maestro Optimizado (PMO) y el periodo regulatorio de la fórmula tarifaria?}
\begin{opts}
    \item El PMO se aplica únicamente a proyectos rurales de Invierte.pe y no tiene relación con la determinación de tarifas.
    \item Ambos instrumentos tienen una vigencia estricta de 12 meses y se recalculan cada vez que varía el tipo de cambio.
    \item El PMO planifica a largo plazo (30 años) y sustenta la fórmula tarifaria quinquenal (5 años).
    \item La fórmula tarifaria fija precios constantes por 30 años sin posibilidad de revisión ni reajuste por inflación.
\end{opts}

\q{Al comparar la regulación por Tasa de Retorno (*Cost of Service*) frente a la regulación por Precios Tope (*Price Cap* con factor de eficiencia $X$), la teoría regulatoria destaca que la regulación por Tasa de Retorno tiende a generar:}
\begin{opts}
    \item Fuertes incentivos para minimizar costos operativos y sobreinvertir en innovación tecnológica ahorradora de capital.
    \item Una eliminación automática de las asimetrías de información entre el regulador y la empresa regulada.
    \item Riesgo financiero excesivo para el operador regulado debido al rezago regulatorio de cinco años.
    \item Sobreinversión en capital (efecto Averch-Johnson) y bajos incentivos para reducir costos.
\end{opts}

\q{En el marco del Sistema Nacional de Programación Multianual y Gestión de Inversiones (Invierte.pe), la diferencia esencial entre la evaluación privada y la evaluación social de un proyecto de saneamiento es que la evaluación social:}
\begin{opts}
    \item Usa precios sociales, incluye externalidades y descuenta a la Tasa Social de Descuento (TSD).
    \item Utiliza precios de mercado con impuestos y descuenta los flujos al WACC privado de la empresa ejecutora.
    \item Excluye los costos de inversión en infraestructura y evalúa exclusivamente la recaudación tarifaria proyectada.
    \item Aplica únicamente a proyectos privados financiados mediante la bolsa de valores sin participación del Estado.
\end{opts}

\q{¿Por qué los impuestos indirectos (como el IGV) y los aranceles incluidos en los precios de mercado se deducen al transformar los costos privados a costos sociales en la evaluación de proyectos de Invierte.pe?}
\begin{opts}
    \item Porque los proyectos públicos de saneamiento están legalmente exonerados de pagar IGV a los contratistas privados.
    \item Porque el Ministerio de Economía y Finanzas asume que la inflación compensa exactamente el efecto del IGV.
    \item Porque son transferencias internas y no representan consumo real de recursos del país.
    \item Porque incluir impuestos indirectos reduciría artificialmente el Valor Actual Neto Privado de la obra.
\end{opts}

\q{En los parámetros de evaluación social publicados por el MEF (Invierte.pe), el Factor de Corrección de la Mano de Obra No Calificada es menor que $1$ (especialmente en zonas rurales), mientras que el Factor de Corrección de la Divisa es mayor que $1$. Esto refleja, respectivamente:}
\begin{opts}
    \item Escasez severa de mano de obra no calificada y sobrevaluación tributaria de las exportaciones peruanas.
    \item Que la evaluación social penaliza el empleo de trabajadores locales y subsidia la importación de maquinaria.
    \item Un error de redondeo estadístico en las cuentas nacionales del Banco Central de Reserva del Perú.
    \item Subempleo estructural (menor costo de oportunidad laboral) y distorsiones que encarecen la divisa.
\end{opts}

\q{En la metodología Costo-Beneficio de Invierte.pe para un proyecto de ampliación de redes de agua potable hacia hogares que antes se abastecían de camiones cisterna, los beneficios sociales directos se cuantifican principalmente mediante:}
\begin{opts}
    \item Ahorro de recursos (cisterna y tiempo) más excedente del consumidor por mayor consumo.
    \item El monto total de impuestos municipales recaudados adicionalmente por la municipalidad provincial.
    \item El costo histórico de construcción de las tuberías multiplicado por la Tasa Social de Descuento.
    \item Exclusivamente la recaudación tarifaria futura que obtendrá la EPS por facturar a los nuevos usuarios.
\end{opts}

\q{En el marco de Invierte.pe para el sector saneamiento, ¿en qué circunstancia corresponde aplicar la metodología de Evaluación Costo-Eficacia (indicador ICE = VAC / unidad de beneficio) en lugar de la metodología Costo-Beneficio (VANS)?}
\begin{opts}
    \item Cuando el proyecto genera un Valor Actual Neto Social negativo y se busca aprobarlo cambiando de indicador.
    \item Cuando los beneficios sociales (ej. salud en PTAR) son difíciles de monetizar con fiabilidad.
    \item Cuando la inversión total del proyecto supera los cien millones de soles y se ejecuta mediante concesión APP.
    \item Cuando todos los insumos del proyecto son importados y no existen factores de corrección del MEF.
\end{opts}

\q{En los lineamientos del MEF para evaluación social en el Perú, los proyectos con horizontes de muy largo plazo orientados a servicios ecosistémicos, conservación de cuencas o mitigación climática (como infraestructura natural vinculada al agua) utilizan tasas sociales de descuento menores que la TSD general ($8\%$) porque:}
\begin{opts}
    \item Los beneficios ecosistémicos siempre se expresan en términos nominales sin ajuste por inflación.
    \item Los proyectos ambientales carecen por completo de costos de operación, mantenimiento y riesgo técnico.
    \item Las normas internacionales prohíben descontar flujos de proyectos hídricos con tasas superiores al $1\%$.
    \item Descontar al $8\%$ casi anula el valor presente de beneficios ambientales futuros.
\end{opts}

\q{En la estructuración de Asociaciones Público-Privadas (APP) de saneamiento en el Perú (como plantas desalinizadoras o PTAR de gran escala), el principio rector de 'Valor por Dinero' (*Value for Money*) establece que se debe elegir la modalidad APP frente a la obra pública tradicional solo si:}
\begin{opts}
    \item Genera mayor calidad o menor costo total ajustado por riesgos durante el ciclo de vida.
    \item Permite ocultar el pasivo contingente fuera del presupuesto fiscal anual sin importar el costo total.
    \item El concesionario privado garantiza que el proyecto será $100\%$ autofinanciado sin ningún cofinanciamiento estatal.
    \item La tasa de financiamiento privado del concesionario es estrictamente inferior a la tasa de los bonos soberanos del Estado.
\end{opts}

\q{De acuerdo con el principio económico de asignación eficiente de riesgos en contratos de concesión APP de saneamiento, cada riesgo del proyecto (construcción, operación, demanda, expropiaciones o fuerza mayor) debe asignarse a:}
\begin{opts}
    \item El concesionario privado en todos los casos, para liberar completamente al Estado de cualquier contingencia fiscal.
    \item La parte (pública o privada) que esté en mejor capacidad de prevenirlo, gestionarlo o mitigarlo al menor costo.
    \item El organismo regulador (SUNASS), que asume la responsabilidad financiera ante sobrecostos de obra.
    \item Los usuarios finales mediante incrementos tarifarios automáticos sin límites ante cualquier ineficiencia constructiva.
\end{opts}

\q{En concesiones cofinanciadas de Plantas de Tratamiento de Aguas Residuales (PTAR) en el Perú, los pagos contractuales RPI (Retribución por Inversiones) y RPMO (Retribución por Operación y Mantenimiento) suelen estructurarse como pagos por disponibilidad condicionados a:}
\begin{opts}
    \item La obtención de utilidades contables positivas por parte de la municipalidad provincial concedente.
    \item El volumen de lluvias anuales registradas en la cuenca hidrográfica sin relación con la operación de la planta.
    \item El cumplimiento de niveles de servicio y calidad del efluente verificado por el supervisor.
    \item La cotización diaria de las acciones del concesionario en la Bolsa de Valores de Lima.
\end{opts}

\q{En el mecanismo peruano de Obras por Impuestos (OxI, Ley N.º 29230) aplicado a proyectos de agua y saneamiento, la empresa privada que financia y ejecuta la obra pública recupera su inversión mediante:}
\begin{opts}
    \item Certificados CIPRL o CIPGN aplicables al pago de su Impuesto a la Renta.
    \item El cobro directo de tarifas reguladas por SUNASS a los usuarios durante un periodo de concesión de treinta años.
    \item La emisión de bonos municipales garantizados por el Fondo de Inversión Agua Segura.
    \item La exoneración permanente del pago del Impuesto General a las Ventas (IGV) en todas sus operaciones comerciales.
\end{opts}

\q{En la evaluación de la sostenibilidad financiera de las EPS del Perú, la 'Relación de Trabajo' (*Working Ratio*, definida como costos operativos sin depreciación divididos entre ingresos operativos) indica una situación financieramente saludable y capacidad de generar flujo para inversiones cuando su valor es:}
\begin{opts}
    \item Estrictamente mayor que $1.00$ (por ejemplo, $1.25$), reflejando alta inversión operativa por cada sol facturado.
    \item Menor que $1.00$ (ej. $0.65$), dejando margen operativo de caja para deuda y CAPEX.
    \item Exactamente igual a $2.00$, lo que asegura la cobertura total de la depreciación y los gastos financieros.
    \item Negativo, lo cual solo ocurre cuando los costos de operación y mantenimiento superan el doble de los ingresos.
\end{opts}

\q{En el esquema regulatorio de SUNASS, los fondos de reserva constituidos a partir de un porcentaje de los ingresos tarifarios para Mecanismos de Retribución por Servicios Ecosistémicos (MERESE) y Gestión de Riesgos de Desastres (GRD) poseen carácter intangible. Esto significa financieramente que:}
\begin{opts}
    \item Pueden utilizarse libremente para pagar planillas administrativas o cubrir déficits de caja corriente de la EPS.
    \item Se transfieren automáticamente al Tesoro Público al cierre de cada ejercicio fiscal si no se ejecutan en el mes.
    \item Deben mantenerse en cuentas separadas solo para su fin aprobado, sin desvío a otros gastos.
    \item Forman parte del patrimonio de libre disposición para distribuir dividendos en efectivo a los municipios accionistas.
\end{opts}

\subsection*{Parte II. Aplicaciones prácticas y ejercicios cuantitativos (154 a 180)}

\q{Un crédito bancario para una EPS pacta una Tasa Nominal Anual (TNA) del $12.0\%$ con capitalización trimestral ($m = 4$). ¿Cuál es la Tasa Efectiva Anual (TEA) equivalente?}
\begin{opts}
    \item $12.00\%$
    \item $12.36\%$
    \item $12.68\%$
    \item $12.55\%$
\end{opts}

\q{Si se invierten $S/\ 1{,}000$ durante $5$ años a una tasa nominal anual del $8\%$ con capitalización continua (considere $e^{0.40} \approx 1.4918$), ¿cuál será el valor futuro acumulado al final del quinto año?}
\begin{opts}
    \item $S/\ 1{,}400.00$
    \item $S/\ 1{,}491.82$
    \item $S/\ 1{,}469.33$
    \item $S/\ 1{,}520.45$
\end{opts}

\q{Si la tasa de descuento nominal de un proyecto de saneamiento es $i = 10.25\%$ anual y la inflación esperada es $\pi = 5.00\%$ anual, ¿cuál es la tasa de descuento real exacta según la ecuación de Fisher?}
\begin{opts}
    \item $4.76\%$
    \item $5.25\%$
    \item $5.00\%$
    \item $5.50\%$
\end{opts}

\q{Calcule el valor presente de una anualidad vencida que paga $S/\ 1{,}000$ al final de cada año durante $3$ años, si la tasa de descuento efectiva anual es $r = 10\%$.}
\begin{opts}
    \item $S/\ 2{,}253.94$
    \item $S/\ 3{,}000.00$
    \item $S/\ 2{,}735.54$
    \item $S/\ 2{,}486.85$
\end{opts}

\q{Una empresa de servicios acaba de pagar un dividendo de $D_0 = S/\ 4.00$ por acción. Si se espera que los dividendos crezcan a perpetuidad a una tasa constante $g = 5\%$ anual y el costo del patrimonio es $k_e = 13\%$, ¿cuál es el precio justo de la acción hoy ($P_0$)?}
\begin{opts}
    \item $S/\ 52.50$
    \item $S/\ 50.00$
    \item $S/\ 48.00$
    \item $S/\ 56.00$
\end{opts}

\q{Una empresa tiene $ROE = 15\%$, reparte el $60\%$ de sus utilidades como dividendos (retiene $b = 40\%$), proyecta una utilidad por acción $UPA_1 = S/\ 10.00$ y enfrenta $k_e = 12\%$. ¿Cuál es su tasa de crecimiento sostenible ($g$) y el precio de su acción ($P_0$)?}
\begin{opts}
    \item $g = 6.0\%$ y $P_0 = S/\ 166.67$
    \item $g = 9.0\%$ y $P_0 = S/\ 200.00$
    \item $g = 6.0\%$ y $P_0 = S/\ 100.00$
    \item $g = 4.0\%$ y $P_0 = S/\ 075.00$
\end{opts}

\q{Al cierre del ejercicio, una EPS presenta Activos Corrientes por $S/\ 450$ millones y Pasivos Corrientes por $S/\ 300$ millones. Calcule el Capital de Trabajo Neto (CTN) y la Razón de Liquidez Corriente.}
\begin{opts}
    \item $CTN = S/\ 750$ millones y Razón Corriente $= 2.50$
    \item $CTN = S/\ 150$ millones y Razón Corriente $= 0.67$
    \item $CTN = S/\ 300$ millones y Razón Corriente $= 1.50$
    \item $CTN = S/\ 150$ millones y Razón Corriente $= 1.50$
\end{opts}

\q{Una empresa registra Activos Corrientes por $S/\ 320$ mil, Inventarios de materiales y repuestos por $S/\ 80$ mil y Pasivos Corrientes por $S/\ 200$ mil. ¿Cuál es su ratio de Prueba Ácida?}
\begin{opts}
    \item $1.20$
    \item $0.80$
    \item $1.60$
    \item $2.00$
\end{opts}

\q{Una empresa concesionaria reporta un margen neto de $8.0\%$, una rotación de activos de $1.25$ veces y un multiplicador del patrimonio ($\text{Activos}/\text{Patrimonio}$) igual a $1.80$. Usando la identidad DuPont, calcule el ROE.}
\begin{opts}
    \item $10.0\%$
    \item $18.0\%$
    \item $14.4\%$
    \item $22.5\%$
\end{opts}

\q{Un proyecto de saneamiento proyecta para el año 1 un $EBIT = S/\ 200$ mil, tasa de impuesto a la renta $T = 30\%$, depreciación de $S/\ 50$ mil, inversiones en activo fijo ($CAPEX$) por $S/\ 80$ mil y un incremento del Capital de Trabajo Neto ($\Delta CTN$) de $S/\ 20$ mil. Calcule el Flujo de Caja Libre ($FCL$).}
\begin{opts}
    \item $S/\ 070$ mil
    \item $S/\ 130$ mil
    \item $S/\ 110$ mil
    \item $S/\ 090$ mil
\end{opts}

\q{Una empresa de servicios públicos registra una Utilidad Operativa Neta después de Impuestos ($NOPAT$) de $S/\ 65$ millones, un Capital Invertido de $S/\ 500$ millones y un $WACC = 10.0\%$. Calcule su Valor Económico Agregado ($EVA$).}
\begin{opts}
    \item $S/\ 15$ millones
    \item $S/\ 10$ millones
    \item $S/\ 50$ millones
    \item $S/\ 65$ millones
\end{opts}

\q{Un proyecto requiere una inversión inicial $I_0 = S/\ 1{,}000$ y genera flujos de caja netos de $S/\ 550$ al final del año 1 y $S/\ 726$ al final del año 2. Si el costo de oportunidad del capital es $r = 10\%$ anual, ¿cuál es el Valor Actual Neto ($VAN$)?}
\begin{opts}
    \item $S/\ 050.00$
    \item $S/\ 100.00$
    \item $S/\ 150.00$
    \item $S/\ 276.00$
\end{opts}

\q{Un proyecto de mejora de redes con inversión inicial $I_0 = S/\ 1{,}000$ miles genera flujos operativos cuyo valor presente descontado al WACC asciende a $S/\ 1{,}320$ miles. ¿Cuál es su Índice de Rentabilidad ($IR$)?}
\begin{opts}
    \item $0.32$
    \item $0.76$
    \item $1.32$
    \item $2.32$
\end{opts}

\q{Determine la Tasa Interna de Retorno ($TIR$) de un proyecto cuya inversión inicial es $I_0 = S/\ 100$ y que produce flujos de caja netos de $S/\ 60$ en el año 1 y $S/\ 72$ en el año 2.}
\begin{opts}
    \item $20.0\%$
    \item $18.0\%$
    \item $15.0\%$
    \item $25.0\%$
\end{opts}

\q{Considere dos proyectos mutuamente excluyentes con igual inversión inicial $I_0 = S/\ 100$: el Proyecto A rinde $S/\ 130$ en $t=1$ y $S/\ 0$ en $t=2$, mientras que el Proyecto B rinde $S/\ 20$ en $t=1$ y $S/\ 121$ en $t=2$. ¿Cuál es la Tasa de Intersección de Fisher ($r_F$)?}
\begin{opts}
    \item $08.0\%$
    \item $10.0\%$
    \item $12.5\%$
    \item $15.0\%$
\end{opts}

\q{Una bomba sumergible tiene un Valor Actual de Costos ($VAC$) de $S/\ 1{,}000$ y una vida útil de $3$ años. Si la tasa de descuento es $r = 10\%$ anual (donde $(1.10)^{-3} \approx 0.7513$), calcule su Costo Anual Equivalente ($CAE$).}
\begin{opts}
    \item $S/\ 333.33$
    \item $S/\ 366.67$
    \item $S/\ 402.11$
    \item $S/\ 433.33$
\end{opts}

\q{Un proyecto requiere $I_0 = S/\ 1{,}000$ y genera flujos netos de $S/\ 550$ en el año 1, $S/\ 484$ en el año 2 y $S/\ 266.20$ en el año 3. Con una tasa de descuento $r = 10\%$ anual, ¿cuál es su Periodo de Recuperación Descontado (*Discounted Payback*)?}
\begin{opts}
    \item $1.93$ años
    \item $2.25$ años
    \item $2.75$ años
    \item $2.50$ años
\end{opts}

\q{Para calcular el costo del capital propio ($k_e$) de una EPS bajo CAPM ajustado, se dispone de una tasa libre de riesgo $r_f = 4.0\%$, un beta apalancado $\beta_L = 0.80$, una prima de riesgo de mercado $(E[R_m] - r_f) = 6.0\%$ y una prima por riesgo país $RP = 2.0\%$. Determine $k_e$.}
\begin{opts}
    \item $08.80\%$
    \item $10.80\%$
    \item $10.00\%$
    \item $12.00\%$
\end{opts}

\q{El beta desapalancado de la industria de saneamiento es $\beta_U = 0.60$. Si una EPS adopta una relación objetivo Deuda/Patrimonio $D/E = 0.50$ y enfrenta una tasa de impuesto a la renta $T = 30\%$, ¿cuál es su beta apalancado ($\beta_L$) según la ecuación de Hamada?}
\begin{opts}
    \item $0.75$
    \item $0.90$
    \item $0.81$
    \item $1.05$
\end{opts}

\q{Una empresa presenta una estructura de capital objetivo de $60\%$ patrimonio y $40\%$ deuda. Si el costo del patrimonio es $k_e = 12.0\%$, el costo de la deuda antes de impuestos es $k_d = 8.0\%$ y la tasa de impuesto a la renta es $T = 25\%$, calcule el $WACC$.}
\begin{opts}
    \item $09.20\%$
    \item $11.00\%$
    \item $10.40\%$
    \item $09.60\%$
\end{opts}

\q{Una empresa financiada íntegramente con capital propio tiene un valor de mercado $V_U = S/\ 800$ millones. Si emite deuda perpetua por $D = S/\ 200$ millones bajo una tasa impositiva corporativa $T = 30\%$ (sin costos de quiebra), ¿cuál será el valor total de la empresa apalancada ($V_L$) según Modigliani-Miller con impuestos?}
\begin{opts}
    \item $S/\ 860$ millones
    \item $S/\ 800$ millones
    \item $S/\ 940$ millones
    \item $S/\ 1{,}000$ millones
\end{opts}

\q{En un mundo sin impuestos (Modigliani-Miller), una empresa tiene un costo de capital desapalancado $k_u = 11.0\%$ y toma deuda a una tasa $k_d = 6.0\%$ hasta alcanzar una razón Deuda/Patrimonio $D/E = 0.80$. ¿Cuál es el nuevo costo del patrimonio ($k_e$)?}
\begin{opts}
    \item $13.5\%$
    \item $16.0\%$
    \item $15.0\%$
    \item $17.0\%$
\end{opts}

\q{Un bono a $2$ años con valor nominal de $S/\ 1{,}000$ paga cupones anuales del $8\%$ ($S/\ 80$ al final de cada año). Si el rendimiento al vencimiento exigido por el mercado es $YTM = 10\%$ anual, ¿cuál es el precio actual del bono?}
\begin{opts}
    \item $S/\ 950.41$
    \item $S/\ 1{,}000.00$
    \item $S/\ 981.82$
    \item $S/\ 965.29$
\end{opts}

\q{Un bono corporativo posee una Duración de Macaulay de $D = 5.40$ años y un rendimiento al vencimiento $y = 8.0\%$ anual. Usando la Duración Modificada, ¿cuál es la variación porcentual aproximada de su precio si el rendimiento sube en $50$ puntos básicos ($\Delta y = +0.50\%$)?}
\begin{opts}
    \item Una caída aproximada de $2.50\%$.
    \item Una caída aproximada de $2.70\%$.
    \item Una caída aproximada de $5.00\%$.
    \item Un alza aproximada de $2.50\%$.
\end{opts}

\q{Una acción que no paga dividendos cotiza a $S_0 = S/\ 100$. Una opción *call* europea a $1$ año con precio de ejercicio $K = S/\ 110.25$ se negocia a $C = S/\ 12.00$. Si la tasa libre de riesgo es $r = 5.0\%$ anual, ¿cuál es el precio de la opción *put* europea ($P$) con igual vencimiento y ejercicio por Paridad Put-Call?}
\begin{opts}
    \item $S/\ 12.00$
    \item $S/\ 17.00$
    \item $S/\ 15.00$
    \item $S/\ 22.25$
\end{opts}

\q{Un fondo obtuvo un retorno $R_p = 14.0\%$ con desviación estándar $\sigma_p = 20.0\%$ y $\beta_p = 1.25$. Si la tasa libre de riesgo es $r_f = 4.0\%$ y el retorno del mercado es $R_m = 10.0\%$, calcule el Alfa de Jensen ($\alpha_p$) y el Ratio de Sharpe ($SR$).}
\begin{opts}
    \item $\alpha_p = +1.5\%$ y $SR = 0.50$
    \item $\alpha_p = +4.0\%$ y $SR = 0.60$
    \item $\alpha_p = +2.5\%$ y $SR = 0.70$
    \item $\alpha_p = +2.5\%$ y $SR = 0.50$
\end{opts}

\q{En la evaluación social bajo Invierte.pe, una obra de agua potable tiene un presupuesto a precios de mercado de $S/\ 1{,}000$ mil, compuesto por $S/\ 600$ mil en insumos nacionales ($FC = 0.85$), $S/\ 200$ mil en mano de obra no calificada ($FC = 0.50$) y $S/\ 200$ mil en equipos importados ($FC = 1.10$). ¿Cuál es el costo de inversión a precios sociales?}
\begin{opts}
    \item $S/\ 830$ mil
    \item $S/\ 810$ mil
    \item $S/\ 780$ mil
    \item $S/\ 890$ mil
\end{opts}
